\section*{Construction eligibility and recorded manual decisions}
\textit{September 11, 2026. Agent-drafted comparison requested by Jacob.}

Jacob asked whether using the same projects for FAR and DUPAC would simplify the
paper without materially changing its results. The proposed rule requires both
density measures to be eligible, finite and positive. Scores, controls, fixed
effects, distance bands and clustering remain as in the current specifications.
This comparison does not adopt the proposed rule in production.

The main comparison would remove 11 multifamily projects from DUPAC, reducing
the full sample from 4,034 to 4,023 and the multifamily sample from 874 to 863.
These projects have usable units and land but unresolved floor area. The main
FAR sample is unchanged. Estimated percentage differences immediately across
the boundary are:

\begin{center}
\begin{tabular}{llrr}
Sample & Outcome & Current & Common sample \\
\hline
All projects & FAR & $-9.24$ & $-9.24$ \\
All projects & DUPAC & $-13.79$ & $-13.70$ \\
Multifamily & FAR & $-12.85$ & $-12.85$ \\
Multifamily & DUPAC & $-17.93$ & $-17.36$ \\
\end{tabular}
\end{center}

The comparison also reproduces the two placebo cutoffs, two donut exclusions,
three boundary-type restrictions and two minimum score gaps. Across these
checks and the main estimates, no coefficient crosses a 1, 5 or 10 percent
significance threshold. All eight placebo estimates remain insignificant at
10 percent. Main and donut estimates remain negative. One FAR observation is
also removed from the positive-placebo sample. Score re-estimation after
own-project permit exclusion is outside this comparison. Table 1 already uses
a citywide common sample and would not change. The recommendation is to use the
common sample for comparability; the cost is discarding valid DUPAC information
for buildings whose floor area remains unresolved.

The current construction Make dependencies reach 26 committed decision tables
containing 809 entries. Entries are not unique projects: they include repeated
project identifiers, source confirmations, exclusions, location decisions and
measurement corrections. This count excludes separate downstream building-type
and zoning reviews. The appendix's 52-project total and single-spreadsheet
description are therefore not a verified description of the current process.
The inventory documents the scope; a count-free replacement describing the
decisions and supporting evidence is recommended but has not been applied.

Jacob separately approved the permit-period correction. Appendix A now states
that permit applications are counted by block and year from 2010 through 2020.
No other manuscript passage or production estimate changes in this step.

\smallskip
\noindent\textit{Verification:} Make builds the 80-row comparison and its
standard report in \texttt{tasks/audits/density\_common\_sample/code/}.
Every estimate, standard error and $p$-value is finite; the common FAR and DUPAC
samples have identical sizes within each specification. A second Make dry run
reports no work. The paper was rebuilt through its Makefile in the completed
replication clone, and the revised appendix page was inspected.
